\documentclass{ctexart}
\usepackage{ctex}
\usepackage{graphicx}
\usepackage{amsmath}
\begin{document}
\title{基于深度学习的中文语义分析研究}
\author{张三}
\keywords{深度学习；语义分析；自然语言处理}
\begin{abstract}
本文提出了一个基于深度学习的中文语义分析方法。实验结果表明，该方法在多个基准数据集上取得了显著的性能提升。
\endabstract}

\section{引言}
近年来，深度学习技术在自然语言处理领域取得了重要进展~\citep{li2024example}。

\section{相关工作}
\begin{equation}
  P(y|\mathbf{x}) = \frac{\exp(\mathbf{w}_y^\top \mathbf{x})}{\sum_{y'} \exp(\mathbf{w}_{y'}^\top \mathbf{x})} \label{eq:softmax}
\end{equation}

\section{方法}
我们提出了一种新的注意力机制。

\section{实验}
\begin{figure}[ht]
  \centering
  \includegraphics[width=0.4\textwidth]{experiment.pdf}
  \caption{实验结果对比。}
  \label{fig:experiment}
\end{figure}

\section{结论}
本文提出了一个最小化的中文学术论文模板，用于测试期刊 Profile 自动检测和 DOCX 生成。

\begin{thebibliography}{99}
  \bibitem{li2024example} 李明. 深度学习在中文语义分析中的应用. \textit{软件学报}, 2024.
\end{thebibliography}
\end{document}
